// Ejercicio 12.1. Longitud de una palabra.
#include <stdio.h>
#include <string.h>

int main() {
   char palabra[100];
   
   printf("Introduce una palabra: ");
   scanf("%s", palabra);
   printf("Longitud: %lu\n", strlen(palabra));
   return 0;
}




// Ejercicio 12.2. Copiar y mostrar.
#include <stdio.h>
#include <string.h>

int main() {
   char original[100];
   char copia[100];
   
   printf("Introduce una palabra: ");
   scanf("%s", original);
   strcpy(copia, original);
   printf("Original: %s\n", original);
   printf("Copia: %s\n", copia);
   return 0;
}



// Ejercicio 12.3. Comparación simple.
#include <stdio.h>
#include <string.h>

int main() {
   char cad1[100], cad2[100];
   
   printf("Introduce la primera palabra: ");
   scanf("%s", cad1);
   
   printf("Introduce la segunda palabra: ");
   scanf("%s", cad2);
   
   if (strcmp(cad1, cad2) == 0) {
      printf("Las palabras son iguales.\n");
   } else {
      printf("Las palabras son diferentes.\n");
   }
   
   return 0;
}


// Ejercicio 12.4. Buscar una letra.
#include <stdio.h>
#include <string.h>

int main() {
   char cadena[100];
   char c;
   
   printf("Introduce una cadena: ");
   scanf("%s", cadena);
   
   printf("Introduce un carácter: ");
   scanf(" %c", &c); // Espacio antes de %c para ignorar el salto de línea
   
   char* pos = strchr(cadena, c);
   
   if (pos != NULL) {
      printf("El carácter '%c' aparece en la posición %ld.\n", c, pos - cadena);
   } else {
      printf("El carácter no se encontró.\n");
   }
   
   return 0;
}


// Ejercicio 12.5. Copia segura con límite.
#include <stdio.h>
#include <string.h>

int main() {
   char origen[100];
   char destino[10];
   
   printf("Introduce una cadena: ");
   scanf("%s", origen);
   
   strncpy(destino, origen, sizeof(destino) - 1);
   destino[sizeof(destino) - 1] = '\0'; // Asegurar el carácter nulo
   
   printf("Destino: %s\n", destino);
   return 0;
}




// Ejercicio 12.6. Comparar prefijos.
#include <stdio.h>
#include <string.h>

int main() {
   char s1[100], s2[100];
   int n;
   
   printf("Introduce la primera cadena: ");
   scanf("%s", s1);
   
   printf("Introduce la segunda cadena: ");
   scanf("%s", s2);
   
   printf("Número de caracteres a comparar: ");
   scanf("%d", &n);
   
   if (strncmp(s1, s2, n) == 0) {
      printf("Coinciden en los primeros %d caracteres.\n", n);
   } else {
      printf("No coinciden en los primeros %d caracteres.\n", n);
   }
   
   return 0;
}





// Ejercicio 12.7. Contar vocales usando strpbrk().
#include <stdio.h>
#include <string.h>

int main() {
   char texto[100];
   const char* vocales = "aeiouAEIOU";
   int contador = 0;
   
   printf("Introduce una cadena: ");
   scanf("%s", texto);
   
   for (int i = 0; texto[i] != '\0'; i++) {
      if (strchr(vocales, texto[i]) != NULL) {
         contador++;
      }
   }
   
   printf("Número de vocales: %d\n", contador);
   return 0;
}




// Ejercicio 12.8. Detectar subcadenas prohibidas.
#include <stdio.h>
#include <string.h>

int main() {
   char frase[200];
   
   printf("Introduce una frase: ");
   getchar(); // Limpiar búfer
   fgets(frase, sizeof(frase), stdin);
   
   if (strstr(frase, "error") != NULL || strstr(frase, "peligro") != NULL ||
         strstr(frase, "prohibido") != NULL) {
      printf("Advertencia: la frase contiene palabras prohibidas.\n");
   } else {
      printf("Frase aceptada.\n");
   }
   
   return 0;
}




// Ejercicio 12.9. Validación numérica.
#include <stdio.h>
#include <string.h>

int main() {
   char codigo[100];
   
   printf("Introduce un código: ");
   scanf("%s", codigo);
   
   if (strspn(codigo, "0123456789") == strlen(codigo)) {
      printf("El código es válido (solo dígitos).\n");
   } else {
      printf("El código contiene caracteres no numéricos.\n");
   }
   
   return 0;
}




// Ejercicio 12.10. Tokenización de una frase.
#include <stdio.h>
#include <string.h>

int main() {
   char frase[200];
   
   printf("Introduce una frase: ");
   fgets(frase, sizeof(frase), stdin);
   
   // Eliminar el salto de línea final
   frase[strcspn(frase, "\n")] = '\0';
   
   char* palabra = strtok(frase, " ");
   
   while (palabra != NULL) {
      printf("Palabra: %s\n", palabra);
      palabra = strtok(NULL, " ");
   }
   
   return 0;
}




// Ejercicio 12.11. Reemplazar todas las letras 'a' por 'x'.
#include <stdio.h>
#include <string.h>

int main() {
   char texto[100];
   
   printf("Introduce una cadena: ");
   scanf("%s", texto);
   
   for (int i = 0; texto[i] != '\0'; i++) {
      if (texto[i] == 'a') {
         texto[i] = 'x';
      }
   }
   
   printf("Cadena modificada: %s\n", texto);
   return 0;
}




// Ejercicio 12.12. Medir la longitud hasta una letra.
   #include <stdio.h>
   #include <string.h>
   
   int main() {
      char texto[100];
      char c;
      
      printf("Introduce una cadena: ");
      scanf("%s", texto);
      
      printf("Introduce el carácter a buscar: ");
      scanf(" %c", &c);
      
      size_t pos = strcspn(texto, (char[]){c, '\0'});
      
      if (pos < strlen(texto)) {
         printf("La letra '%c' aparece tras %zu caracteres.\n", c, pos);
      } else {
         printf("El carácter no se encuentra en la cadena.\n");
      }
      
      return 0;
   }

